\section{Supplementary Proof Appendix}
% 补充证明附录
\label{app:proofs}

This appendix is clearly marked as supplementary material and does not enlarge the model or claims of the main paper.
% 本附录明确标记为补充材料，不扩大正文的模型或主张。
It supplies the proof objects used by \cref{thm:exact-history} for one finite registered protected-call slice and one atomic policy-domain epoch at a time.
% 它给出 \cref{thm:exact-history} 在每次一个有限已注册受保护调用切片和一个原子策略域纪元范围内所使用的证明对象。
Independent domains retain the product interpretation stated in the paper.
% 独立域仍采用正文所述的乘积解释。

\subsection{The Declarative Specification and Its Computed Witness}
% 声明式规范及其计算见证
\label{app:spec-bridge}

Fix a well-formed admission instance \(\Theta=(\kappa,H,\Delta,\mathsf{out},\mathcal A,u)\) for which the history judgment derives \(\delta=(H^+,\mathcal C,\eta^-,\eta^+)\).
% 固定一个良构的准入实例 \(\Theta=(\kappa,H,\Delta,\mathsf{out},\mathcal A,u)\)，其历史判定派生出 \(\delta=(H^+,\mathcal C,\eta^-,\eta^+)\)。
All sets in this subsection are indexed sets, so two equal resolved words belonging to different outcomes remain different elements.
% 本小节中的所有集合都是带索引集合，因此属于不同结果的两个相同已解析词仍是不同元素。
\begin{definition}[Source-promise coverage]
% 源承诺覆盖
\label{app:spc}
\[
\begin{aligned}
\mathsf{SPC}(H,u,\widehat B)
&\Longleftrightarrow
\forall a\in\mathsf{Req}_u.\ \exists t\in O_{H,u},b.\\[-2pt]
&\hspace{32mm}\mathsf{Lift}_u(t,a)\land(t,b)\in\widehat B .
\end{aligned}
\]
Here \(\mathsf{Dec}_{H,u}\) is the typed losing-arm outcomes for \(\mathsf{MergeSelect}\), or the current-role outcomes of a pending certified \(\mathsf{RestoreReplace}\); \(\mathsf{Done}_{H,u}\) is the completed current role of \(\mathsf{RestoreReplace}\); both are empty for every other row and for live extension.
% 此处 \(\mathsf{Dec}_{H,u}\) 是 \(\mathsf{MergeSelect}\) 的类型化失败分支结果，或待处理且经认证的 \(\mathsf{RestoreReplace}\) 的当前角色结果；\(\mathsf{Done}_{H,u}\) 是 \(\mathsf{RestoreReplace}\) 已完成的当前角色；对于其余各行与在线扩展，二者均为空。
Thus \(\mathsf{Req}_u=\mathsf{SrcProm}(H,u)\setminus(\mathsf{Done}_{H,u}\uplus\mathsf{Dec}_{H,u})\), so \(\mathsf{SPC}\) covers exactly the promises neither already discharged nor explicitly declassified.
% 因此 \(\mathsf{Req}_u=\mathsf{SrcProm}(H,u)\setminus(\mathsf{Done}_{H,u}\uplus\mathsf{Dec}_{H,u})\)，所以 \(\mathsf{SPC}\) 恰好覆盖既未履行也未被显式解除的承诺。
\end{definition}
Let \(\mathfrak R_\Theta\) be the finite set of indexed realizations from \cref{def:realization}, ordered componentwise by
% 令 \(\mathfrak R_\Theta\) 为 \cref{def:realization} 中有限的带索引实现集合，并按如下分量关系排序：
\[
 (L_1,\widehat B_1)\sqsubseteq(L_2,\widehat B_2)
 \quad\Longleftrightarrow\quad
 L_1\subseteq L_2\ \land\ \widehat B_1\subseteq\widehat B_2 .
\]

\begin{lemma}[Declarative--algorithmic bridge]
% 声明式--算法桥接
\label{lem:app-spec-bridge}
The declarative predicate \(\mathsf{GReal}(\Theta;L,\widehat B)\) has a witness if and only if \(\widehat B^\dagger_{H,u}\ne\varnothing\).
% 当且仅当 \(\widehat B^\dagger_{H,u}\ne\varnothing\) 时，声明式谓词 \(\mathsf{GReal}(\Theta;L,\widehat B)\) 才有见证。
When a witness exists, it is unique and equals \((W^\dagger_{H,u},\widehat B^\dagger_{H,u})\).
% 当见证存在时，它是唯一的，并且等于 \((W^\dagger_{H,u},\widehat B^\dagger_{H,u})\)。
This equality is a theorem relating two independently stated objects, rather than the definition of the ideal machine.
% 这一等式是联系两个独立陈述对象的定理，而不是理想机的定义。
\end{lemma}

\begin{proof}
Let \((L,\widehat B)\in\mathfrak R_\Theta\).
% 令 \((L,\widehat B)\in\mathfrak R_\Theta\)。
Source and target fidelity and policy safety give \(\widehat B\subseteq\widehat B_0\).
% 源与目标忠实性以及策略安全性给出 \(\widehat B\subseteq\widehat B_0\)。
For every \((o,b)\in\widehat B\), every prefix \(z\preceq_{\rm pref}b\) belongs to \(L=\Pref(\pi_2\widehat B)\).
% 对每个 \((o,b)\in\widehat B\)，每个前缀 \(z\preceq_{\rm pref}b\) 都属于 \(L=\Pref(\pi_2\widehat B)\)。
Persistent outcome coverage therefore supplies, for every \(v\in\mathsf{Compat}(z)\), a pair \((v,b')\in\widehat B\) extending \(z\).
% 因而，持续结果覆盖为每个 \(v\in\mathsf{Compat}(z)\) 提供一个扩展 \(z\) 的二元组 \((v,b')\in\widehat B\)。
Thus \(\widehat B\subseteq\widehat\Phi(\widehat B)\).
% 因此 \(\widehat B\subseteq\widehat\Phi(\widehat B)\)。
By monotonicity on the finite lattice \(2^{\widehat B_0}\), every post-fixed set is contained in the greatest fixed point, so \(\widehat B\subseteq\widehat B^\dagger\).
% 由于有限格 \(2^{\widehat B_0}\) 上的单调性，每个后不动点集合都包含于最大不动点，因此 \(\widehat B\subseteq\widehat B^\dagger\)。
Taking projected prefixes yields \(L\subseteq W^\dagger\).
% 对投影取前缀可得 \(L\subseteq W^\dagger\)。

Conversely, suppose \(\widehat B^\dagger\ne\varnothing\).
% 反之，假设 \(\widehat B^\dagger\ne\varnothing\)。
The fixed-point equality supplies persistent coverage and, at \(\epsilon\), composes with checked lifts to establish \(\mathsf{SPC}\); inclusion in \(\widehat B_0\) supplies target fidelity and safe completed words, prefix closure of \(\mathcal A\) supplies safety for every generated prefix, and \(W^\dagger=\Pref(\pi_2\widehat B^\dagger)\) supplies completion nonblocking.
% 不动点等式提供持续覆盖，并在 \(\epsilon\) 处与经检查的提升组合以建立 \(\mathsf{SPC}\)；包含于 \(\widehat B_0\) 提供目标忠实性与安全完成词，\(\mathcal A\) 的前缀闭包为每个生成前缀提供安全性，而 \(W^\dagger=\Pref(\pi_2\widehat B^\dagger)\) 提供完成非阻塞性。
Hence \((W^\dagger,\widehat B^\dagger)\in\mathfrak R_\Theta\), and the preceding containment makes it greatest.
% 因此 \((W^\dagger,\widehat B^\dagger)\in\mathfrak R_\Theta\)，并且前述包含关系使其成为最大元素。
If two greatest elements existed, each would contain the other componentwise, so they would be equal.
% 如果存在两个最大元素，则它们按分量彼此包含，因此必然相等。
If \(\widehat B^\dagger=\varnothing\), the first half places every hypothetical realization inside the empty set, contradicting its required nonemptiness.
% 如果 \(\widehat B^\dagger=\varnothing\)，证明的前半部分会把每个假想实现包含于空集，这与实现必须非空相矛盾。
\end{proof}

\begin{corollary}[Exact compilation boundary]
% 精确编译边界
\label{cor:app-exact-boundary}
The ideal history edge is enabled by the declarative \(\mathsf{Spec}(\Theta)\) exactly when the checked compiler returns a nonempty fixed point.
% 理想历史边由声明式 \(\mathsf{Spec}(\Theta)\) 启用，当且仅当经检查的编译器返回非空不动点。
On that edge, the ideal language \(L^*\) equals the generated language of \(\mathsf{Can}(W^\dagger,B^\dagger)\).
% 在该边上，理想语言 \(L^*\) 等于 \(\mathsf{Can}(W^\dagger,B^\dagger)\) 的生成语言。
\end{corollary}

\begin{proof}
The first statement is \cref{lem:app-spec-bridge}.
% 第一项陈述即为 \cref{lem:app-spec-bridge}。
The second combines that lemma with \cref{lem:compiled-language}.
% 第二项陈述由该引理与 \cref{lem:compiled-language} 共同得到。
\end{proof}

\begin{proposition}[Output-sensitive compilation]
% 输出敏感编译
\label{prop:app-output-complexity}
Fix the unit-cost symbolic-checker model with constant-time finite-map lookup and policy-automaton transitions.
% 固定单位代价的符号检查器模型，其中有限映射查询与策略自动机转移均为常数时间。
Let
\[
\begin{aligned}
D&=\lvert\Theta\rvert_{\rm chk},&
n&=\max_o\lvert X_{p_o}\rvert,\\
\Lambda&=\sum_o\sum_{w\in\Lin(p_o)}(\lvert w\rvert+1),&
K&=\lvert\widehat B_0\rvert,\\
L&=\sum_{(o,b)\in\widehat B_0}(\lvert b\rvert+1).
\end{aligned}
\]
\[
\Gamma=\{((o,b),z,v)\mid
(o,b)\in\widehat B_0,\ z\preceq b,\ v\in\mathsf{Compat}(z)\}.
\]
Here \(D\) is the authenticated structural input size, \(\Lambda\) is indexed linearization volume, and \(\Gamma\) is the materialized prefix--outcome obligation relation.
% 其中 \(D\) 是经认证结构输入大小，\(\Lambda\) 是带索引线性化体积，而 \(\Gamma\) 是物化的前缀--结果义务关系。
The semantic compiler runs in \(O(D+n\Lambda+|\Gamma|+L\log(2+L))\) time and \(O(D+\Lambda+|\Gamma|)\) space, including its emitted witnesses.
% 语义编译器在 \(O(D+n\Lambda+|\Gamma|+L\log(2+L))\) 时间和 \(O(D+\Lambda+|\Gamma|)\) 空间内运行，其中包括所输出的见证。
Moreover, \(K\leq\sum_o|\Lin(p_o)|\), \(|\Gamma|\leq K(n+1)|O_{\mathcal C_u}|\), and the number of nonempty strict-removal ranks is at most \(K\).
% 此外，\(K\leq\sum_o|\Lin(p_o)|\)、\(|\Gamma|\leq K(n+1)|O_{\mathcal C_u}|\)，且非空严格移除秩的数量至多为 \(K\)。
\end{proposition}

\begin{proof}
A backtracking topological-order enumerator emits all indexed linearizations in \(O(n\Lambda)\) time, while resolution, policy simulation, and construction of raw and resolved prefix tries are linear in emitted volume.
% 回溯式拓扑序枚举器在 \(O(n\Lambda)\) 时间内发出全部带索引线性化，而解析、策略模拟以及原始与已解析前缀字典树的构造关于输出体积均为线性。
For every support key \((z,v)\), maintain the number of live pairs \((v,b')\) extending \(z\), initially enqueue candidates having a zero-count obligation, and remove candidates in synchronous batches.
% 对每个支撑键 \((z,v)\)，维护扩展 \(z\) 的存活二元组 \((v,b')\) 数量，初始时将具有零计数义务的候选入队，并按同步批次删除候选。
When removal makes a support count zero, enqueue exactly the candidates incident to that obligation for the next batch.
% 当删除使一个支撑计数变为零时，恰将关联该义务的候选放入下一批次。
Each candidate and each incidence in \(\Gamma\) is processed at most once, and the batches are exactly \(\widehat B_i\setminus\widehat B_{i+1}\).
% 每个候选以及 \(\Gamma\) 中的每个关联至多处理一次，并且这些批次恰为 \(\widehat B_i\setminus\widehat B_{i+1}\)。
Store one rank and cause per removed candidate rather than copied sets, and construct \(\mathsf{Can}(W,B)\) by bottom-up sorting of finite-trie derivative signatures in \(O(L\log(2+L))\) time.
% 对每个被删除候选只存储一项秩与原因而不复制集合，并通过自底向上排序有限字典树的导数签名，在 \(O(L\log(2+L))\) 时间内构造 \(\mathsf{Can}(W,B)\)。
\end{proof}

\begin{proof}[Proof of \cref{lem:live-extension}]
% \cref{lem:live-extension} 的证明
\label{app:live-extension-proof}
An authenticated extension is identity-on-frontier, preserves cell resolution and all indexed outcomes, and satisfies \(\mathcal A\subseteq\mathcal A^+\), so every old safe completion and coverage witness remains valid for the successor operator.
% 经认证扩展在前沿上保持恒等，保留单元解析与全部带索引结果，并满足 \(\mathcal A\subseteq\mathcal A^+\)，因此每条旧的安全完成路径与覆盖见证对后继算子仍然有效。
The current residual family is therefore post-fixed, and greatestness gives its inclusion in a nonempty successor fixed point.
% 因而，当前残余族是后不动点，而最大性给出它包含于非空的后继不动点中。
\end{proof}

\subsection{Lean Theorem Map}
% Lean 定理映射
\label{app:lean-map}

\path{RegistrationRefinement.checkRegistration_iff}, \path{RegistrationRefinement.compiled_cell_eq_iff_canonical_eq}, \path{RegistrationRefinement.compiled_mediation_exhaustive}, the three receipt theorems, and \path{RegistrationRefinement.compiled_admission_preserved} cover registration reflection and its link to admission.
% \path{RegistrationRefinement.checkRegistration_iff}、\path{RegistrationRefinement.compiled_cell_eq_iff_canonical_eq}、\path{RegistrationRefinement.compiled_mediation_exhaustive}、三条回执定理与 \path{RegistrationRefinement.compiled_admission_preserved} 覆盖注册反射及其与准入的连接。
\path{FiniteCore.compiler_eq_semantic} and the greatest-realization theorems cover the fixed point, while \path{HistoryStructure.deriveEdit_iff} covers all six structural derivations.
% \path{FiniteCore.compiler_eq_semantic} 与最大实现定理覆盖不动点，而 \path{HistoryStructure.deriveEdit_iff} 覆盖全部六种结构推导。
\path{CompilationBridge.compiler_admit_exists_secure_install} constructs a valid installed cut from admission, and its converse declarations recover admission from any valid install candidate.
% \path{CompilationBridge.compiler_admit_exists_secure_install} 从准入构造有效安装切点，而其逆向声明从任意有效安装候选恢复准入。
\path{OperationalSemantics.kernelTrace_preserves_agentSec}, \path{OperationalSemantics.live_extension_preserves_agentSec}, \path{OperationalSemantics.fresh_before_live_extension_stale}, \path{OperationalSemantics.alias_before_live_extension_stale}, and \path{OperationalSemantics.live_extension_before_old_use_denied} cover installed finite-prefix closure and the extension races.
% \path{OperationalSemantics.kernelTrace_preserves_agentSec}、\path{OperationalSemantics.live_extension_preserves_agentSec}、\path{OperationalSemantics.fresh_before_live_extension_stale}、\path{OperationalSemantics.alias_before_live_extension_stale} 与 \path{OperationalSemantics.live_extension_before_old_use_denied} 覆盖安装后的有限前缀闭包与扩展竞态。
The audit module imports these declarations and prints only the project-allowed axioms \texttt{propext}, \texttt{Classical.choice}, and \texttt{Quot.sound}.
% 审计模块导入这些声明，并且只打印项目允许的公理 \texttt{propext}、\texttt{Classical.choice} 与 \texttt{Quot.sound}。

\subsection{Normalized Evidence and Safe Projection}
% 规范化证据与安全投影
\label{app:evidence}

We make the normalization implicit in \(\mathsf{Sec}\) explicit.
% 我们将 \(\mathsf{Sec}\) 中隐含的规范化过程显式化。
Let \(\mathcal Q_{\rm sec}\) pair every valid \(V=\mathsf{Sec}(S;c,r)\) with a well-sorted \(q::=\mathsf{Decide}(u)\mid\mathsf{TryInstall}(u,Y)\), permitting stale names.
% 令 \(\mathcal Q_{\rm sec}\) 将每个有效 \(V=\mathsf{Sec}(S;c,r)\) 与允许过期名称的良好定型 \(q::=\mathsf{Decide}(u)\mid\mathsf{TryInstall}(u,Y)\) 配对。
\(\mathsf{Ans}\) returns \(\mathsf{Compile}(\Theta_V(u))\) for \(\mathsf{Decide}\), and for \(\mathsf{TryInstall}\) returns the successor bundle exactly when \(\mathsf{Ready}_V(u,Y)\), otherwise \(\mathsf{NoCommit}\).
% 对 \(\mathsf{Decide}\)，\(\mathsf{Ans}\) 返回 \(\mathsf{Compile}(\Theta_V(u))\)；对 \(\mathsf{TryInstall}\)，它恰在 \(\mathsf{Ready}_V(u,Y)\) 时返回后继包，否则返回 \(\mathsf{NoCommit}\)。
An equivariant \(\tau\) is exact iff some \(D_\tau\) maps every \(\langle\tau(V),q\rangle_{\cong}\) to \([\mathsf{Ans}(V,q)]_{\cong}\).
% 等变 \(\tau\) 是精确的，当且仅当存在某个 \(D_\tau\) 将每个 \(\langle\tau(V),q\rangle_{\cong}\) 映射到 \([\mathsf{Ans}(V,q)]_{\cong}\)。
Let \(\mathcal K\) be the finite typed key universe at a checked cut, containing tagged outcome, occurrence, cell, receipt, checkpoint, group, domain, epoch, policy-version, and schema-version keys.
% 令 \(\mathcal K\) 为一个经检查切点处的有限类型化键宇宙，包含带标签的结果、发生项、单元、回执、检查点、组、域、纪元、策略版本和模式版本键。
Public keys are authority labels and identifiers declared external by the interface, while every other key may be consistently renamed.
% 公开键包括权限标签和接口声明为外部的标识符，而其他每个键都可以一致地重命名。
For a typed query \(q\), \(\mathsf{supp}(q)\) is defined recursively over its full syntax, including all identifiers nested inside a supplied symbolic certificate, monitor, or cut seal.
% 对带类型查询 \(q\)，\(\mathsf{supp}(q)\) 在其完整语法上递归定义，包括所提供符号证书、监控器或切点封印中嵌套的全部标识符。
Let \(\mathcal F\) be the finite set of typed field paths and assign every semantic base field a unique owner
% 令 \(\mathcal F\) 为有限的类型化字段路径集合，并为每个语义基础字段指定唯一所有者：
\[
 \mathsf{own}:\mathcal F_{\rm base}
 \longrightarrow\{\mathbf P,\mathbf I,\mathbf E,\mathbf C\}.
\]
\(\mathbf P\) owns history constructor tags, live/checkpoint/target contract nodes, order and outcome membership, schema kind and source fingerprints, lift relations, and retirement declarations.
% \(\mathbf P\) 拥有历史构造器标签、活跃/检查点/目标契约节点、顺序与结果成员关系、模式种类与源指纹、提升关系和退役声明。
\(\mathbf I\) owns occurrence--cell incidence, cell--canonical-invocation/label/scope rows, occurrence lineage, and logical gate and handle slots.
% \(\mathbf I\) 拥有发生项--单元关联、单元--规范调用/标签/作用域记录、发生项谱系以及逻辑门与句柄槽位。
\(\mathbf E\) owns the global resolved trace, live and checkpoint cursors, receipt--cell/creator and retry paths, result futures, ordered outbox phase rows, and residual policy.
% \(\mathbf E\) 拥有全局已解析轨迹、活跃与检查点游标、回执--单元/创建者与重试路径、结果期值、有序发件箱阶段记录和残余策略。
\(\mathbf C\) owns the domain, policy/schema/view versions, private namespace, epoch allocation and phase, active epoch and gate ownership, certificate origin and coordinates, and checkpoint cut coordinates.
% \(\mathbf C\) 拥有域、策略/模式/视图版本、私有命名空间、纪元分配与阶段、活跃纪元与门所有权、证书来源与坐标以及检查点切点坐标。
These four lists are disjoint and exhaustive for semantic base fields.
% 这四组列表对语义基础字段既互不相交又完全覆盖。
Certificate bytes, seals, fingerprints, digests, derived authority words, and foreign-key paths are derived fields rather than additional semantic facts.
% 证书字节、封印、指纹、摘要、派生权限词和外键路径是派生字段，而不是额外的语义事实。

Write \(f\leadsto g\) when the value or well-typed interpretation of derived field \(g\) depends immediately on field \(f\).
% 当派生字段 \(g\) 的值或良类型解释直接依赖字段 \(f\) 时，记 \(f\leadsto g\)。
The dependency graph is acyclic because canonical encodings are constructed from base relations in a fixed order.
% 依赖图是无环的，因为规范编码按固定顺序从基础关系构造。
Every foreign-key path depends on its target; authenticated gate and handle bytes depend on their \(\mathbf I\) slot and \(\mathbf C\) epoch/owner; receipt reconstruction and the authority word depend on \(\mathbf E\) receipt/cursor rows and their \(\mathbf I\) cells and labels; and target certificates, cut seals, and digests depend on every encoded base field in their payload.
% 每条外键路径依赖其目标；认证门与句柄字节依赖其 \(\mathbf I\) 槽位和 \(\mathbf C\) 纪元/所有者；回执重建与权限词依赖 \(\mathbf E\) 的回执/游标记录及其 \(\mathbf I\) 单元与标签；目标证书、切点封印和摘要则依赖其载荷中编码的每个基础字段。
A normalized full view \(V\) is a compatible finite assignment to all base fields together with the unique values of all derivable fields.
% 规范化完整视图 \(V\) 是对所有基础字段的相容有限赋值，并带有全部可派生字段的唯一值。
Compatibility requires equal values on shared typed keys and the well-formedness conditions of \cref{def:well-formed-history}.
% 相容性要求共享类型化键取值相同，并满足 \cref{def:well-formed-history} 的良构条件。
The natural join \(\mathbf P\bowtie\mathbf I\bowtie\mathbf E\bowtie\mathbf C\) is the compatible union of these uniquely owned base fields followed by deterministic recomputation of derived fields.
% 自然连接 \(\mathbf P\bowtie\mathbf I\bowtie\mathbf E\bowtie\mathbf C\) 是这些具有唯一所有权的基础字段的相容并集，随后确定性地重新计算派生字段。
It is undefined precisely when shared typed keys disagree or a required foreign-key target is absent.
% 当且仅当共享类型化键不一致或必需的外键目标缺失时，该连接无定义。

For \(J\subseteq\{\mathbf P,\mathbf I,\mathbf E,\mathbf C\}\), start with base fields owned by \(J\) and the public sort declarations needed to type the retained records.
% 对 \(J\subseteq\{\mathbf P,\mathbf I,\mathbf E,\mathbf C\}\)，先保留由 \(J\) 所有的基础字段和为保留记录定型所需的公开类型声明。
Repeatedly delete any foreign-key path whose target has been deleted and any derived field having a deleted transitive dependency.
% 随后反复删除目标已被删除的任何外键路径，以及具有已删除传递依赖的任何派生字段。
The unique fixed point of this deletion is \(\mathsf{Ret}(J)\), and \(\pi_J(V)\) is \(V\) restricted to \(\mathsf{Ret}(J)\).
% 这一删除过程的唯一不动点记为 \(\mathsf{Ret}(J)\)，而 \(\pi_J(V)\) 是 \(V\) 在 \(\mathsf{Ret}(J)\) 上的限制。
This definition prevents an erased fact from surviving inside a seal, digest, archived receipt-only registry path, or certificate payload.
% 该定义防止被擦除事实继续隐藏在封印、摘要、仅供回执使用的归档注册表路径或证书载荷中。

The incidence-erasure projection \(\pi_{\neg oc}\) starts from the full view, deletes the occurrence--cell relation, and applies the same dependency deletion while retaining the occurrence and cell marginals.
% 关联擦除投影 \(\pi_{\neg oc}\) 从完整视图出发，删除发生项--单元关系，并应用相同的依赖删除，同时保留发生项和单元边际记录。
The projection \(\pi_{\neg rc}\) analogously deletes receipt--cell incidence and every path that reconstructs it, including receipt--creator, cursor--receipt, and archived creator--cell paths, while retaining the receipt, creator, occurrence, and cell marginals.
% 投影 \(\pi_{\neg rc}\) 类似地删除回执--单元关联以及每条可重建该关联的路径，包括回执--创建者、游标--回执和归档创建者--单元路径，同时保留回执、创建者、发生项和单元边际记录。
State views are compared modulo a sort-preserving bijection on private keys that fixes all public keys.
% 状态视图在保持类型的私有键双射意义下比较，该双射固定全部公开键。
For a common literal query \(q\), write \(V\cong_qV'\) when the witnessing bijection additionally fixes \(\mathsf{supp}(q)\) pointwise; equivalently, \(\langle V,q\rangle_{\cong}=\langle V',q\rangle_{\cong}\).
% 对共同的字面查询 \(q\)，当见证双射还逐点固定 \(\mathsf{supp}(q)\) 时，记 \(V\cong_qV'\)；等价地，\(\langle V,q\rangle_{\cong}=\langle V',q\rangle_{\cong}\)。

\begin{lemma}[Projection is well defined]
% 投影良定义性
\label{lem:app-projection}
Natural join, dependency deletion, and private renaming commute.
% 自然连接、依赖删除和私有重命名彼此可交换。
Consequently, \([\pi_J(V)]_{\cong}\), \([\pi_{\neg oc}(V)]_{\cong}\), and \([\pi_{\neg rc}(V)]_{\cong}\) do not depend on a representative or on a certificate serialization.
% 因此，\([\pi_J(V)]_{\cong}\)、\([\pi_{\neg oc}(V)]_{\cong}\) 和 \([\pi_{\neg rc}(V)]_{\cong}\) 均不依赖某个代表元或证书序列化方式。
\end{lemma}

\begin{proof}
A sort-preserving bijection preserves key equality, foreign-key reachability, ownership, and the dependency relation.
% 保持类型的双射会保持键相等性、外键可达性、所有权和依赖关系。
It therefore maps each deletion round to the corresponding deletion round in the renamed view.
% 因而，它把每轮删除映射到重命名视图中的对应删除轮次。
Induction on the finite acyclic dependency graph proves that the retained fixed points correspond.
% 对有限无环依赖图作归纳可证明保留的不动点彼此对应。
Canonical recomposition is equivariant because it is a typed constructor over exactly those retained fields.
% 规范重组是等变的，因为它是恰好作用于这些保留字段的类型化构造器。
\end{proof}

\begin{lemma}[Representation-independent observational separation]
% 与表示无关的观测分离
\label{lem:app-observational-separation}
Let \(V_0,V_1\) be finite valid full views and let \(q\) be a typed query.
% 令 \(V_0,V_1\) 为有限有效完整视图，并令 \(q\) 为带类型查询。
If \([\mathsf{Ans}(V_0,q)]_{\cong}\ne[\mathsf{Ans}(V_1,q)]_{\cong}\), then every exact abstraction \(\tau\) satisfies \(\tau(V_0)\not\cong_q\tau(V_1)\).
% 若 \([\mathsf{Ans}(V_0,q)]_{\cong}\ne[\mathsf{Ans}(V_1,q)]_{\cong}\)，则每个精确抽象 \(\tau\) 都满足 \(\tau(V_0)\not\cong_q\tau(V_1)\)。
The claim applies to arbitrary encodings and is not a minimal-table or unique-representation claim.
% 该结论适用于任意编码，并不声称表的数量最少或表示方式唯一。
\end{lemma}

\begin{proof}
Assume toward contradiction that \(\tau(V_0)\cong_q\tau(V_1)\).
% 反设 \(\tau(V_0)\cong_q\tau(V_1)\)。
The joint decoder inputs \(\langle\tau(V_0),q\rangle_{\cong}\) and \(\langle\tau(V_1),q\rangle_{\cong}\) then coincide.
% 此时，联合解码器输入 \(\langle\tau(V_0),q\rangle_{\cong}\) 与 \(\langle\tau(V_1),q\rangle_{\cong}\) 相同。
Exactness forces equal answer orbits, contradicting the premise.
% 精确性迫使答案轨道相等，与前提矛盾。
\end{proof}

\begin{table}[t]
\centering
\scriptsize
\setlength{\tabcolsep}{2pt}
\begin{tabular}{@{}c >{\raggedright\arraybackslash}p{0.47\columnwidth}
                  >{\raggedright\arraybackslash}p{0.34\columnwidth}@{}}
\toprule
Erased & Full-view pair & Exact answers\\
\midrule
\multicolumn{3}{@{}l}{\(u_R(\mathcal C):=\mathsf{RestoreReplace}(b,k,\sigma)\) with \(\mathcal C_k=\mathcal C\).}\\
\(\mathbf P\) &
\(q_{\mathbf P}=\mathsf{MergeSelect}(g,L,\sigma)\);
\(T^0=x\oplus_g^{\rm open}y\), \(T^1=x\parallel_g y\); otherwise equal. &
\(V^0\): realize; \(V^1\): \(\mathsf{Invalid}\) (wrong mode).\\
\(\mathbf I\) &
\(q_{\mathbf I}=\mathsf{RestoreLive}(b_L,k,\sigma_{\mathbf I})\) after safe empty-root/ForkChoice/Save;
\(\Delta=\epsilon,\mathcal A=\Pref\{a\}\);
\(q^0_{\rm cell}(x)=q^0_{\rm cell}(y)=d_0\);
\(q^1_{\rm cell}(x)=d_0\ne d_1=q^1_{\rm cell}(y)\);
\(\ell(d_0)=\ell(d_1)=a\). &
\(X\parallel Y_k\): one \Fresh admit;
two \Fresh reject \(aa\).\\
\(\mathbf E\) &
\(q_{\mathbf E}=\mathsf{TryInstall}(u_R(x),Y_0)\);
\(Y_0\) is compiled at \(V^0\) after safe empty-root/ForkChoice/Save/\Fresh;
\(V^0\xrightarrow{\Dispatch(d)}V^1\), changing only the release phase. &
\(V^0\): install;
\(V^1\): \(\mathsf{NoCommit}\) (stale cut).\\
\(\pi_{\neg rc}\) &
\(q_{\neg rc}=u_R(x)\) after safe empty-root/ForkChoice/Save/\Fresh;
\(q_{\rm cell}(x)=d\), \(\mathcal A=\Pref\{a\}\);
same marginals \(\{p\},\{d,e\}\);
\(p\mapsto d\) in \(V^0\), \(p\mapsto e\) in \(V^1\). &
\(\Alias(x_k,d)\): admit;
\(\Fresh(x_k,d)\): reject \(aa\).\\
\(\mathbf C\) &
\(q_{\mathbf C}=\mathsf{TryInstall}(u_0,Y_0)\), with full package \(Y_0\) compiled at \(V^0\), is literal in both;
name supplies agree off the epoch sort, while live epoch, ownership, and dependent seals differ. &
\(V^0\): install; \(V^1\): reject stale cut.\\
\bottomrule
\end{tabular}
\caption{Factor- and incidence-erasure pairs with isomorphic retained views and opposite exact answers.}
% 因子与关联擦除对具有同构的保留视图和相反的精确答案。
\label{tab:projection-pairs}
\end{table}

\subsection{Bootstrap States and Explicit Separation Pairs}
% 引导状态与显式分离对
\label{app:boot-pairs}

We construct the finite lower-bound witnesses from one literal empty genesis \(\mathsf S_\emptyset\).
% 我们从唯一的字面空起源状态 \(\mathsf S_\emptyset\) 构造有限下界见证。
A registered slice \(R\) compiled from a boundary-finite typed Agent workflow \(\mathcal W\) contains only finite contracts, schemas, policy, lowering data, and authenticated occurrence, cell, invocation, lineage, scope, and label rows.
% 从边界有限的类型化 Agent 工作流 \(\mathcal W\) 编译得到的注册切片 \(R\) 只包含有限契约、模式、策略、降低数据，以及经认证的发生项、单元、调用、谱系、作用域和标签行。
It contains no receipt, cursor, outbox, epoch, certificate, result future, monitor state, or installed cut.
% 它不包含回执、游标、发件箱、纪元、证书、结果期值、监控器状态或已安装切点。
The pre-slice relation has a registration step \(\mathsf S_\emptyset\xrightarrow{\mathsf{register}(\mathcal W,R,n)}\mathsf{Reg}(R,n)\) exactly when \(\mathsf{CheckReg}(\mathcal W,R)\) accepts and \(n\) is a fresh private namespace, followed by checked root installation.
% 当且仅当 \(\mathsf{CheckReg}(\mathcal W,R)\) 接受且 \(n\) 是新鲜私有命名空间时，切片前关系才包含注册步骤 \(\mathsf S_\emptyset\xrightarrow{\mathsf{register}(\mathcal W,R,n)}\mathsf{Reg}(R,n)\)，随后执行经检查的根安装。
Root installation derives the initial cut and constructs empty receipt, cursor, and outbox state, generation zero, one active epoch, the canonical monitor, and authenticated live bindings.
% 根安装推导初始切点，并构造空回执、空游标与空发件箱状态、零代、一个活跃纪元、规范监控器和经认证的活跃绑定。

\begin{lemma}[Registration reflection]
% 注册反射
\label{lem:registration-reflection}
For finite-state \(\mathcal W\) and finite \(R\), \(\mathsf{CheckReg}(\mathcal W,R)\) accepts iff the registered \(\lambda\) witnesses \(\mathcal W\sqsubseteq_\lambda R\) in \cref{def:registration-refinement}.
% 对于有限状态的 \(\mathcal W\) 与有限的 \(R\)，当且仅当已注册的 \(\lambda\) 见证 \cref{def:registration-refinement} 中的 \(\mathcal W\sqsubseteq_\lambda R\) 时，\(\mathsf{CheckReg}(\mathcal W,R)\) 才接受。
On acceptance, the \(\lambda\)-lowerings of source executions modulo \(\equiv_\partial\) correspond bijectively to \(\mathsf{BRuns}(R)\).
% 接受时，模 \(\equiv_\partial\) 的源执行经 \(\lambda\) 降低后与 \(\mathsf{BRuns}(R)\) 双射对应。
The correspondence preserves indexed outcomes, causal order, canonical invocation quotienting, receipt evolution, and the admission answer.
% 该对应保持带索引结果、因果次序、规范调用商、回执演化与准入答案。
\end{lemma}
\begin{proof}
The productive-SCC test terminates because the state graph is finite.
% 有效强连通分量检查因状态图有限而终止。
A productive SCC containing a durable-site edge can be pumped before an outcome, producing projected words of unbounded length.
% 含有持久位置边的有效强连通分量可以在到达结果前被反复遍历，从而产生长度无界的投影词。
Conversely, if no productive SCC contains such an edge, every productive SCC is boundary-silent, and contracting them yields a DAG with a finite enumerable projected language.
% 反之，若任何有效强连通分量都不含此类边，则每个有效强连通分量在边界上静默，收缩它们会得到具有有限可枚举投影语言的 DAG。
The checker compares that language with \(\mathsf{BRuns}(R)\) and directly decides outcome/site incidence, order, cell attestation, total unique lowering, raw-bypass exclusion, and prefix receipt commutation.
% 检查器将该语言与 \(\mathsf{BRuns}(R)\) 比较，并直接判定结果/位置关联关系、次序、单元证明、全定义唯一降低、原始绕过排除与前缀回执交换性。
Therefore it accepts exactly the semantic conjunction defining \(\mathcal W\sqsubseteq_\lambda R\), rather than an oracle restatement.
% 因而，它恰好接受定义 \(\mathcal W\sqsubseteq_\lambda R\) 的语义合取，而不是对预言机的重述。
Boundary projection erases only internal steps, so matching projected paths induce the quotient-execution bijection.
% 边界投影只擦除内部步骤，因此匹配的投影路径诱导执行商上的双射。
Canonical invocation equality and deterministic \(\Resolve\) make receipt evolution commute with this bijection at every prefix.
% 规范调用相等性与确定性的 \(\Resolve\) 使回执演化在每个前缀处与该双射交换。
The realization families are therefore order-isomorphic, so the fixed-point operator preserves admission, greatestness, and pullback of the monitor.
% 因而，实现族彼此序同构，所以不动点算子保持准入、最大性与监控器回拉。
\end{proof}

\begin{lemma}[Natural bootstrap]
% 自然引导
\label{lem:app-boot}
If \(\mathsf{CheckReg}(\mathcal W,R)\) accepts and \(R\)'s registered root contract and policy have a nonempty greatest realization, registration and checked root installation reach a finite state \(S_R\) with \(\mathsf{AgentSec}(S_R;c_R,\epsilon)\).
% 若 \(\mathsf{CheckReg}(\mathcal W,R)\) 接受且 \(R\) 的已注册根契约与策略具有非空的最大实现，则注册与经检查的根安装会到达一个满足 \(\mathsf{AgentSec}(S_R;c_R,\epsilon)\) 的有限状态 \(S_R\)。
The cut origin is the checked root derivation, while every later cut-changing transition records a six-edit derivation or an authenticated live extension.
% 该切点的来源是经检查的根推导，而此后每个改变切点的转移都记录一种六类编辑推导或一次经认证的在线扩展。
\end{lemma}

\begin{proof}
Registration contributes only the checked \(R\) and \(n\), and deterministic allocation constructs every runtime identifier from them.
% 注册只提供经检查的 \(R\) 和 \(n\)，而确定性分配据此构造每个运行时标识。
The root judgment and checked lifts establish core/schema coherence, and the recomputed nonempty fixed point establishes promise coherence.
% 根判定与经检查的提升建立核心/模式一致性，而重新计算的非空不动点建立承诺一致性。
Empty runtime progress establishes execution coherence, the canonical program at \(q_\epsilon\) establishes monitor coherence, and the atomic phase/binding assignments establish epoch coherence.
% 空运行时进度建立执行一致性，位于 \(q_\epsilon\) 的规范程序建立监控器一致性，而原子的阶段/绑定赋值建立纪元一致性。
Only initial installation creates a root origin, and every successful later installation stores its checked structural-edit or live-extension derivation.
% 只有初始安装会创建根来源，而此后每次成功安装都会存储其经检查的结构编辑或在线扩展推导。
\end{proof}

The following finite pairs use private names except for each displayed query and authority labels.
% 下列有限历史对除所显示的查询和权限标签外均使用私有名称。
Each displayed state is reached from \(\mathsf S_\emptyset\) by the stated registration, root-installation, Save, protected-use, and edit steps; omitted fields are identical up to the stated private renaming and dependency deletion.
% 每个显示状态都从 \(\mathsf S_\emptyset\) 经由所述注册、根安装、Save、受保护使用和编辑步骤到达；省略字段在所述私有重命名与依赖删除意义下相同。

\paragraph{\(\mathbf P\): structural topology}
% \(\mathbf P\)：结构拓扑
Let \(X\) and \(Y\) be singleton contracts with distinct occurrences \(x\) and \(y\), common cells, and a universal residual policy.
% 令 \(X\) 和 \(Y\) 为具有不同发生项 \(x\) 与 \(y\)、公共单元和全称残余策略的单元素契约。
From the same registered singleton root, checked initial installation followed by \(\mathsf{ForkChoice}\) reaches \(V_{\mathbf P}^0\), while the corresponding \(\mathsf{ForkParallel}\) trace reaches \(V_{\mathbf P}^1\).
% 从同一个已注册单元素根出发，经检查的初始安装随后执行 \(\mathsf{ForkChoice}\) 会到达 \(V_{\mathbf P}^0\)，而对应的 \(\mathsf{ForkParallel}\) 轨迹会到达 \(V_{\mathbf P}^1\)。
The two views differ only in
% 两个视图仅在下式中不同：
\[
 V_{\mathbf P}^0.T=b_L{:}X\oplus_g^{\mathsf{open}}b_R{:}Y,
 \qquad
 V_{\mathbf P}^1.T=b_L{:}X\parallel_g b_R{:}Y .
\]
Use the same registered \(\mathsf{MergeSelect}\) schema and query \(q_{\mathbf P}=\mathsf{MergeSelect}(g,L,\sigma)\).
% 使用相同的已注册 \(\mathsf{MergeSelect}\) 模式和查询 \(q_{\mathbf P}=\mathsf{MergeSelect}(g,L,\sigma)\)。
The choice view has a unique structural derivation and a nonempty one-step greatest realization, whereas the parallel view has no \(\mathsf{MergeSelect}\) derivation because its group mode is wrong.
% 选择视图具有唯一的结构推导和非空的一步最大实现，而并行视图由于组模式错误而没有 \(\mathsf{MergeSelect}\) 推导。
Deleting \(\mathbf P\) removes the constructor tag and every authenticator depending on it, leaving isomorphic retained views.
% 删除 \(\mathbf P\) 会移除构造器标签及依赖它的每个认证器，从而留下同构的保留视图。

\paragraph{Bootstrap convention for the incidence pairs}
% 关联对的引导约定
Let \(\mathbf 1\) be the singleton-outcome contract whose unique pomset is empty.
% 令 \(\mathbf 1\) 为唯一偏序事件集为空的单结果契约。
Its completion \(\epsilon\) gives a nonempty greatest realization under \(\mathcal A=\Pref\{a\}\).
% 其完成路径 \(\epsilon\) 在 \(\mathcal A=\Pref\{a\}\) 下给出非空最大实现。
Both pairs below install this safe root and then a checked \(\mathsf{ForkChoice}\).
% 下述两对都安装这一安全根，再安装一次经检查的 \(\mathsf{ForkChoice}\)。
Each dormant schema row needed to equalize cell marginals is identical in both views and never becomes live.
% 为使单元边际记录相等而需要的每条休眠模式记录在两个视图中都相同且永不激活。

\paragraph{\(\mathbf I\): occurrence--cell incidence}
% \(\mathbf I\)：发生项--单元关联
Let \(X,Y\) be singleton contracts with occurrences \(x,y\), and let \(d_0,d_1\) be unreceipted cells labeled \(a\).
% 令 \(X,Y\) 为分别具有发生项 \(x,y\) 的单元素契约，并令 \(d_0,d_1\) 为标签为 \(a\) 的无回执单元。
Two checked registered slices have identical topology, schemas, occurrence marginals, and cell marginals but the following incidence:
% 两个经检查的注册切片具有相同的拓扑、模式、发生项边际记录与单元边际记录，但具有如下关联关系：
\[
\begin{array}{c|cc}
 &q_{\rm cell}(x)&q_{\rm cell}(y)\\ \hline
 V_{\mathbf I}^0&d_0&d_0\\
 V_{\mathbf I}^1&d_0&d_1 .
\end{array}
\]
From the installed \(\mathbf 1\) root, the common Fork schema derives \(b_L{:}X\oplus_g^{\mathsf{open}}b_R{:}Y\).
% 从已安装的 \(\mathbf 1\) 根出发，公共 Fork 模式派生 \(b_L{:}X\oplus_g^{\mathsf{open}}b_R{:}Y\)。
Each outcome redeems one cell, so both Forks install under \(\Pref\{a\}\).
% 每个结果只兑付一个单元，所以两次 Fork 都在 \(\Pref\{a\}\) 下安装。
Save \(b_R\) before progress and issue the common query \(q_{\mathbf I}=\mathsf{RestoreLive}(b_L,k,\sigma_{\mathbf I})\).
% 在取得进展前保存 \(b_R\)，再发出公共查询 \(q_{\mathbf I}=\mathsf{RestoreLive}(b_L,k,\sigma_{\mathbf I})\)。
Up to clone renaming, the target is \((X\otimes Y_k)\oplus Y\).
% 在克隆重命名意义下，目标为 \((X\otimes Y_k)\oplus Y\)。
In \(V_{\mathbf I}^0\), the left outcome uses one \Fresh and one \Alias on \(d_0\), while the bypass outcome uses one \Fresh, so every outcome has authority word \(a\).
% 在 \(V_{\mathbf I}^0\) 中，左侧结果在 \(d_0\) 上使用一次 \Fresh 与一次 \Alias，而绕过结果使用一次 \Fresh，所以每个结果的权限词都是 \(a\)。
In \(V_{\mathbf I}^1\), every left-outcome linearization uses \Fresh on both \(d_0,d_1\), so its word is \(aa\notin\mathcal A\).
% 在 \(V_{\mathbf I}^1\) 中，左侧结果的每条线性化都在 \(d_0,d_1\) 上使用 \Fresh，所以其权限词为 \(aa\notin\mathcal A\)。
The left outcome is compatible with \(\epsilon\), so the first \(\widehat\Phi\) round also removes the otherwise safe bypass completion and leaves the greatest fixed point empty.
% 左侧结果与 \(\epsilon\) 兼容，因此第一轮 \(\widehat\Phi\) 也会移除本来安全的绕过完成路径，并使最大不动点为空。
Erasing occurrence--cell incidence and its dependent rows makes the retained views isomorphic while preserving equal marginals.
% 擦除发生项--单元关联及其依赖记录会使保留视图同构，同时保持相等的边际记录。

\paragraph{\(\mathbf E\): mutable effect progress}
% \(\mathbf E\)：可变效果进度
Let \(C,X\) be singleton contracts with occurrences \(c,x\), let \(q_{\rm cell}(c)=q_{\rm cell}(x)=d\), and label \(d\) by \(a\).
% 令 \(C,X\) 为分别具有发生项 \(c,x\) 的单元素契约，令 \(q_{\rm cell}(c)=q_{\rm cell}(x)=d\)，并以 \(a\) 标记 \(d\)。
From the safe \(\mathbf 1\) root, a checked Fork installs \(b_L{:}C\oplus_g^{\mathsf{open}}b_R{:}X\); save \(b_R\), then use \(c\) to complete the left arm, create receipt \(p\), and prepare its outbox item.
% 从安全的 \(\mathbf 1\) 根出发，一次经检查的 Fork 安装 \(b_L{:}C\oplus_g^{\mathsf{open}}b_R{:}X\)；保存 \(b_R\)，再使用 \(c\) 完成左侧分支、创建回执 \(p\) 并准备其发件箱条目。
At the resulting \(V_{\mathbf E}^0\), compile \(u_{\mathbf E}=\mathsf{RestoreReplace}(b_L,k,\sigma_{\mathbf E})\) to the full package \(Y_0\); its clone of \(x\) aliases \(d\), so \(Y_0\) is installable under \(\mathcal A=\Pref\{a\}\).
% 在所得 \(V_{\mathbf E}^0\) 处，将 \(u_{\mathbf E}=\mathsf{RestoreReplace}(b_L,k,\sigma_{\mathbf E})\) 编译为完整包 \(Y_0\)；其 \(x\) 的克隆会别名到 \(d\)，所以 \(Y_0\) 在 \(\mathcal A=\Pref\{a\}\) 下可安装。
Let \(V_{\mathbf E}^0\xrightarrow{\Dispatch(d)}V_{\mathbf E}^1\), which changes only \(\mathbf E\)'s outbox-release phase and preserves \(\mathbf P,\mathbf I,\mathbf C\) literally.
% 令 \(V_{\mathbf E}^0\xrightarrow{\Dispatch(d)}V_{\mathbf E}^1\)；该步骤仅改变 \(\mathbf E\) 的发件箱释放阶段，并逐字保持 \(\mathbf P,\mathbf I,\mathbf C\)。
For the common literal \(q_{\mathbf E}=\mathsf{TryInstall}(u_{\mathbf E},Y_0)\), \(V_{\mathbf E}^0\) installs, whereas \(V_{\mathbf E}^1\) returns \(\mathsf{NoCommit}\) because the cut seal binds the outbox phase.
% 对于共同字面查询 \(q_{\mathbf E}=\mathsf{TryInstall}(u_{\mathbf E},Y_0)\)，\(V_{\mathbf E}^0\) 完成安装，而 \(V_{\mathbf E}^1\) 返回 \(\mathsf{NoCommit}\)，因为切点封印绑定了发件箱阶段。

\paragraph{\(\pi_{\neg rc}\): receipt--cell incidence}
% \(\pi_{\neg rc}\)：回执--单元关联
For a separate pair, retain \(q_{\rm cell}(x)=d\) and cells \(d,e\) labeled \(a\), but let the archived creator bind \(p\mapsto d\) in \(V_{\neg rc}^0\) and \(p\mapsto e\) in \(V_{\neg rc}^1\).
% 对于另一对状态，保持 \(q_{\rm cell}(x)=d\) 以及以 \(a\) 标记的单元 \(d,e\)，但令归档创建者在 \(V_{\neg rc}^0\) 中绑定 \(p\mapsto d\)，在 \(V_{\neg rc}^1\) 中绑定 \(p\mapsto e\)。
Both states follow the same safe empty-root/ForkChoice/Save/\Fresh shape above, have authority word \(a\), and use the common query \(q_{\neg rc}=\mathsf{RestoreReplace}(b_L,k,\sigma_{\neg rc})\).
% 两个状态都遵循上述安全的空根/ForkChoice/Save/\Fresh 形状，具有权限词 \(a\)，并使用共同查询 \(q_{\neg rc}=\mathsf{RestoreReplace}(b_L,k,\sigma_{\neg rc})\)。
The cloned \(x_k\) is \(\Alias(x_k,d)\) in \(V_{\neg rc}^0\), but \(\Fresh(x_k,d)\) produces forbidden word \(aa\) in \(V_{\neg rc}^1\).
% 克隆的 \(x_k\) 在 \(V_{\neg rc}^0\) 中为 \(\Alias(x_k,d)\)，而在 \(V_{\neg rc}^1\) 中，\(\Fresh(x_k,d)\) 会产生被禁止的权限词 \(aa\)。
Here the only differing occurrence--cell row belongs to the archived receipt creator; \(\pi_{\neg rc}\) explicitly deletes that reconstructing path together with receipt--cell incidence, leaving isomorphic retained views with equal receipt, occurrence, and cell marginals.
% 此处唯一不同的发生项--单元记录属于已归档的回执创建者；\(\pi_{\neg rc}\) 会连同回执--单元关联显式删除这条重建路径，从而留下具有相等回执、发生项和单元边际记录的同构保留视图。

\paragraph{\(\mathbf C\): history cut}
% \(\mathbf C\)：历史切点
Let a bootstrap namespace be a family \(n=(n_s)_{s\in\mathsf{Sort}}\) of injective, per-sort fresh supplies.
% 令引导命名空间为一族 \(n=(n_s)_{s\in\mathsf{Sort}}\)；其中每个类型的供给均为单射且新鲜。
From the same literal empty genesis, register the same public slice \(R\) with supplies \(n^0,n^1\) that agree on every non-epoch sort but mint distinct initial epochs \(\eta_0\ne\eta_1\), and perform checked root installation.
% 从同一个字面空起源出发，使用供给 \(n^0,n^1\) 注册同一个公开切片 \(R\)；二者在每个非纪元类型上相同，但生成不同的初始纪元 \(\eta_0\ne\eta_1\)，随后执行经检查根安装。
The resulting \(V_{\mathbf C}^0,V_{\mathbf C}^1\) have literally identical \(\mathbf P,\mathbf I,\mathbf E\) base fields, while \(\mathbf C\)'s active epoch, gate ownership, certificate coordinates, and dependent authenticators differ.
% 得到的 \(V_{\mathbf C}^0,V_{\mathbf C}^1\) 具有字面完全相同的 \(\mathbf P,\mathbf I,\mathbf E\) 基础字段，而 \(\mathbf C\) 的活跃纪元、门所有权、证书坐标及其依赖认证器不同。
Compile the same registered edit \(u_0\) at \(V_{\mathbf C}^0\) to obtain \(Y_0=\mathsf{Installable}(\delta_0,W_0,\widehat B_0,Z_0)\), and submit the same literal query \(q_{\mathbf C}=\mathsf{TryInstall}(u_0,Y_0)\) to both views.
% 在 \(V_{\mathbf C}^0\) 处编译同一个已注册编辑 \(u_0\) 得到 \(Y_0=\mathsf{Installable}(\delta_0,W_0,\widehat B_0,Z_0)\)，并向两个视图提交同一个字面查询 \(q_{\mathbf C}=\mathsf{TryInstall}(u_0,Y_0)\)。
At \(V_{\mathbf C}^0\), recomputation reproduces \(\mathsf{cut}_{Z_0}\), so installation succeeds.
% 在 \(V_{\mathbf C}^0\) 处，重新计算会重现 \(\mathsf{cut}_{Z_0}\)，因此安装成功。
At \(V_{\mathbf C}^1\), recomputation binds \(\eta_1\) and its owner rather than the \(\eta_0\)-coordinates sealed in \(Z_0\), so installation rejects the stale cut.
% 在 \(V_{\mathbf C}^1\) 处，重新计算绑定 \(\eta_1\) 及其所有者，而非 \(Z_0\) 中封印的 \(\eta_0\) 坐标，因此安装会拒绝这一过期切点。
Erasing \(\mathbf C\) recursively erases epoch-dependent gate/handle rows, certificate coordinates, seals, and authenticators; the identity bijection on every retained sort then witnesses \(V_{\mathbf C}^0\cong_{q_{\mathbf C}}V_{\mathbf C}^1\).
% 擦除 \(\mathbf C\) 会递归擦除依赖纪元的门/句柄记录、证书坐标、封印和认证器；此后每个保留类型上的恒等双射见证 \(V_{\mathbf C}^0\cong_{q_{\mathbf C}}V_{\mathbf C}^1\)。

\begin{proposition}[Explicit factor and incidence lower bounds]
% 显式因子与关联下界
\label{prop:app-explicit-pairs}
Every pair above consists of finite states reached from the common empty genesis and satisfying \(\mathsf{AgentSec}\).
% 上述每一对均由从共同空起源状态到达且满足 \(\mathsf{AgentSec}\) 的有限状态组成。
Each pair has \(q\)-isomorphic indicated projections and opposite answer orbits for its common literal query.
% 每一对都具有关于 \(q\) 同构的指定投影，并对其共同字面查询给出相反的答案轨道。
Therefore every proper factor projection, \(\pi_{\neg oc}\), and \(\pi_{\neg rc}\) is inexact.
% 因此，每个真因子投影、\(\pi_{\neg oc}\) 和 \(\pi_{\neg rc}\) 都是不精确的。
Every exact abstraction must distinguish all displayed pairs.
% 每个精确抽象都必须区分所有显示的配对。
\end{proposition}

\begin{proof}
\Cref{lem:app-boot} gives the common genesis and valid root installation, while the explicit traces above use only modeled Save, protected-use, and six-edit transitions.
% \Cref{lem:app-boot} 给出共同起源状态和有效根安装，而上述显式轨迹只使用已建模的 Save、受保护使用和六种编辑转移。
For the incidence pairs, the empty root completes with \(\epsilon\), and each installed Fork arm has a one-\Fresh completion labeled \(a\).
% 对于关联对，空根以 \(\epsilon\) 完成，并且每个已安装 Fork 分支都具有一条标签为 \(a\) 的单次 \Fresh 完成路径。
Hence every displayed pre-query state is reachable under \(\Pref\{a\}\).
% 因此，每个显示的查询前状态在 \(\Pref\{a\}\) 下均可达。
Open choice enables the Save steps, and the installed monitors enable each stated left-arm \Fresh before that arm becomes a completed, still-addressable continuation.
% 开放选择启用这些 Save 步骤，而已安装监控器会在左侧分支成为已完成但仍可寻址的延续之前启用各个所述左侧分支 \Fresh。
For the \(\mathbf E\) pair, \Dispatch preserves \(\mathbf P,\mathbf I,\mathbf C\) but changes an outbox phase bound by \(Y_0\)'s cut seal, giving install versus \(\mathsf{NoCommit}\); the displayed schema and resolution calculations give the other opposite answers.
% 对于 \(\mathbf E\) 配对，\Dispatch 保持 \(\mathbf P,\mathbf I,\mathbf C\)，但改变 \(Y_0\) 的切点封印所绑定的发件箱阶段，从而分别得到安装与 \(\mathsf{NoCommit}\)；所显示的模式与解析计算给出其余相反答案。
\Cref{lem:app-projection} and the specified private-name bijections, each pointwise fixing the common query support, give \(q\)-relative projection equality.
% \Cref{lem:app-projection} 与所指定的私有名称双射给出关于 \(q\) 的投影相等性；每个双射都逐点固定共同查询的支撑。
Any projection omitting at least one factor collapses the corresponding factor pair, while the two incidence erasures collapse their corresponding incidence pairs.
% 任何遗漏至少一个因子的投影都会合并对应的因子对，而两种关联擦除会合并各自对应的关联对。
Inexactness and the general information lower bound now follow from \cref{lem:app-observational-separation}.
% 不精确性与一般信息下界由 \cref{lem:app-observational-separation} 得出。
\end{proof}

Assigning the same caller-proposed target plant, markings, and epoch to the \(\mathbf P\) pair makes \(\tau_{\rm tgt}\) equal while leaving the authenticated queries and opposite answers unchanged.
% 为 \(\mathbf P\) 对赋予相同的由调用方提出的目标被控对象、标记和纪元，会使 \(\tau_{\rm tgt}\) 相等，同时保持经认证的查询及其相反答案不变。
Thus target-only pre-derivation state is not exact, although generalized nonblocking is exact as a backend after trusted derivation supplies the plant, outcome identities, and conditional markings.
% 因而，仅含目标的派生前状态并不精确，尽管可信派生提供被控对象、结果标识与条件标记之后，广义非阻塞作为后端是精确的。

\subsection{Event Derivatives and Invariant Preservation}
% 事件导数与不变量保持
\label{app:event-preservation}

Write \(A_1,\ldots,A_5\) for the five clauses of \(\mathsf{AgentSec}\) in their stated order.
% 依照陈述顺序，用 \(A_1,\ldots,A_5\) 表示 \(\mathsf{AgentSec}\) 的五个条款。
The following matrix is exhaustive for durable transitions inside an installed registered protected-call slice.
% 下表穷尽一个已安装的已注册受保护调用切片内部的持久转移。
Initial registration occurs before root admission; later additive schema and policy growth is the authenticated live-extension history edge listed below.
% 初始注册发生在根准入之前；后续仅追加的模式与策略增长由下表所列的经认证在线扩展历史边完成。

\begin{table*}[t]
\centering
\scriptsize
\setlength{\tabcolsep}{3pt}
\begin{tabular}{@{}p{0.15\textwidth}p{0.29\textwidth}p{0.12\textwidth}p{0.36\textwidth}@{}}
\toprule
Event class & Durable update & Witness & Preservation argument\\
\midrule
\Fresh use &
Apply \(\mathsf{HStep}\) to \((G,T)\); append \(\Fresh(x,d)\) to global \(\chi\); advance \(\nu,q\); append one receipt and canonical-invocation outbox item; remove live binding. &
\((c,r\Fresh(x,d))\) &
\(A_1\): registry and schema unchanged.
\(A_2\): residual coherence.
\(A_3\): unique receipt and ordered preparation.
\(A_4\): canonical derivative.
\(A_5\): remove exactly \(x\)'s binding.\\
\Alias use &
Apply \(\mathsf{HStep}\) to \((G,T)\); append \(\Alias(x,d)\) to global \(\chi\); advance \(\nu,q\); add only the result-future link; remove live binding. &
\((c,r\Alias(x,d))\) &
\(A_1\): authenticated cell retained.
\(A_2\): residual coherence.
\(A_3\): link names an earlier receipt and authority is unchanged.
\(A_4\): canonical derivative.
\(A_5\): remove exactly \(x\)'s binding.\\
Successful history edge, six edits or live extension &
Install derived \((\kappa^+,H^+,\mathcal A^+)\), checked monitor and certificate; close \(\eta^-\); activate \(\eta^+\); publish new bindings. &
\((c_\delta^+,\epsilon)\) &
\(A_1\): schema fidelity and lifts.
\(A_2\): declarative--algorithmic bridge and checker.
\(A_3\): source receipts and outbox retained.
\(A_4\): install at \(q_\epsilon\).
\(A_5\): total domain handoff.\\
Verified \(\mathsf{Invalid}\) or \(\mathsf{Reject}\) &
Emit the canonical current-cut certificate; no durable state mutation. &
\((c,r)\) &
All clauses unchanged; \(\gamma\) binds the visible answer to this decision cut.\\
\(\mathsf{Save}\) &
Snapshot the current residual/cursor pair in an authenticated checkpoint; advance \(\nu\). &
\((c,r)\) &
\(A_1\): immutable well-typed extension.
\(A_2,A_4,A_5\): unchanged.
\(A_3\): exactly the permitted \(K\) extension.\\
\(\mathsf{Preload}\) &
Allocate or replace content only in an inactive unbound epoch. &
\((c,r)\) &
\(A_1\)--\(A_4\): unchanged.
\(A_5\): inactive and unbound remain true; no handle is exposed.\\
Retry, retrieval, delivery &
Emit \(\mathsf{Return}(t,[v]_{\cong})\) for an authenticated durable value; no security-state mutation. &
\((c,r)\) &
All clauses unchanged; retry additionally checks the same canonical invocation and receipt link.\\
Denial &
Return \(\mathsf{deny}(x)\); no state mutation. &
\((c,r)\) &
All clauses unchanged.\\
\(\Dispatch(d)\) &
Mark exactly the oldest prepared outbox item released. &
\((c,r)\) &
\(A_3\): release order remains a prefix of preparation order.
Other clauses unchanged.\\
Settlement &
Emit \(\mathsf{Settle}(d,[v]_{\cong})\), advance one receipt phase, and fill its stable result future. &
\((c,r)\) &
\(A_3\): cell, invocation, creator, authority, and order are immutable.
Other clauses unchanged.\\
Stale or malformed candidate &
No durable state mutation. &
\((c,r)\) &
All clauses unchanged; the attempt is \(\tau\) and emits no verdict.\\
Crash and recovery &
Discard volatile work and expose the same recovered durable image. &
\((c,r)\) &
All clauses unchanged under the assumed crash-stable transaction boundary.\\
\bottomrule
\end{tabular}
\caption{Exhaustive event classes and preservation of the five \(\mathsf{AgentSec}\) clauses.}
% 穷尽的事件类别及 \(\mathsf{AgentSec}\) 五个条款的保持。
\label{tab:app-event-matrix}
\end{table*}

\begin{lemma}[Representative-independent event derivative]
% 与代表元无关的事件导数
\label{lem:app-event-derivative}
The guarded update induces a partial function \(\partial\) on joint state--event orbits \(\langle V,e\rangle_{\cong}\).
% 带守卫更新在联合状态--事件轨道 \(\langle V,e\rangle_{\cong}\) 上诱导一个偏函数 \(\partial\)。
If \(\mathsf{AgentSec}(S;c,r)\) and \(S\xrightarrow eS'\), the witness in \cref{tab:app-event-matrix} satisfies \(\mathsf{AgentSec}(S';c',r')\) and
% 若 \(\mathsf{AgentSec}(S;c,r)\) 且 \(S\xrightarrow eS'\)，则 \cref{tab:app-event-matrix} 中的见证满足 \(\mathsf{AgentSec}(S';c',r')\)，并且
\[
 [\mathsf{Sec}(S';c',r')]_{\cong}
 =\partial(\langle\mathsf{Sec}(S;c,r),e\rangle_{\cong}).
\]
\end{lemma}

\begin{proof}
All rule guards use typed equality, membership, order, canonical constructors, and authenticated payload equality.
% 所有规则守卫仅使用类型化相等性、成员关系、次序、规范构造器和经认证载荷相等性。
These operations are equivariant under a simultaneous private-key renaming of the state and event payload.
% 这些操作在状态与事件载荷的同步私有键重命名下是等变的。
For every such renaming \(\pi\), \(G_e(V)\Longleftrightarrow G_{\pi e}(\pi V)\) and \(\mathsf U_{\pi e}(\pi V)=\pi\mathsf U_e(V)\).
% 对每个这样的重命名 \(\pi\)，都有 \(G_e(V)\Longleftrightarrow G_{\pi e}(\pi V)\) 且 \(\mathsf U_{\pi e}(\pi V)=\pi\mathsf U_e(V)\)。
Thus equal pointed orbits have the same definedness and the same successor orbit.
% 因而，相等的带指向轨道具有相同的有定义性和相同的后继轨道。
This proves representative independence and functionality.
% 这证明了与代表元无关性和函数性。

For \Fresh and \Alias, \cref{lem:resolution,lem:residual-coherence,lem:compiled-language} establish the first two rows of the matrix.
% 对于 \Fresh 与 \Alias，\cref{lem:resolution,lem:residual-coherence,lem:compiled-language} 建立表中的前两行。
If \(x\) lies at \(b\), \(\mathsf{HStep}\) appends \(a\) to \(b\)'s cursor, while the same atomic protected-use update appends \(a\) to the global resolved-occurrence trace \(\chi\); residual associativity gives
% 若 \(x\) 位于 \(b\)，\(\mathsf{HStep}\) 会把 \(a\) 追加到 \(b\) 的游标，而同一个原子受保护使用更新会把 \(a\) 追加到全局已解析发生项轨迹 \(\chi\)；残余结合律给出
\[
 (\Gamma_b/\mathsf{raw}(\chi_b))/x
 =\Gamma_b/\mathsf{raw}(\chi_ba).
\]
Thus the updated leaf and any later Save remain synchronized with \(G\).
% 因此，更新后的叶节点以及此后的任何 Save 都与 \(G\) 保持同步。
For a successful history edge, \cref{lem:history-derivation,lem:live-extension,lem:app-spec-bridge,lem:compiled-language,lem:cut-linearization} establish the third row for all six schema rows and authenticated live extension.
% 对于成功历史边，\cref{lem:history-derivation,lem:live-extension,lem:app-spec-bridge,lem:compiled-language,lem:cut-linearization} 为六个模式行与经认证在线扩展建立表中的第三行。
Each remaining rule changes only the fields shown in the matrix, and its row checks every invariant clause whose fields change.
% 其余每条规则只改变表中显示的字段，而对应行会检查字段发生变化的每个不变量条款。
Thus every successor has the displayed witness and normalized update.
% 因此，每个后继都具有所显示的见证与规范化更新。
\end{proof}

\begin{corollary}[Finite-prefix closure]
% 有限前缀闭包
\label{cor:app-prefix-closure}
Starting from any valid bootstrap, every finite prefix of any sequence of the event classes in \cref{tab:app-event-matrix} ends in a state satisfying \(\mathsf{AgentSec}\).
% 从任意有效引导状态出发，\cref{tab:app-event-matrix} 中事件类别所组成任意序列的每个有限前缀都终止于满足 \(\mathsf{AgentSec}\) 的状态。
The sequence length is unbounded, although every individual admission instance and event payload is finite.
% 序列长度不受固定上界限制，但每个单独的准入实例与事件载荷都是有限的。
For a finite word \(\bar e\), the iterated derivative is defined on \(\langle V,\bar e\rangle_{\cong}\), where one renaming acts on the initial view and the whole word.
% 对于有限词 \(\bar e\)，迭代导数定义在 \(\langle V,\bar e\rangle_{\cong}\) 上，其中同一次重命名作用于初始视图与整个事件词。
\end{corollary}

\begin{proof}
The base case is \cref{lem:app-boot}, and \cref{lem:app-event-derivative} supplies both the invariant-preserving inductive step and equivariance under one renaming of the remaining event suffix.
% 基础情形由 \cref{lem:app-boot} 给出，而 \cref{lem:app-event-derivative} 同时给出保持不变量的归纳步骤，以及在同一次重命名下对剩余事件后缀的等变性。
\end{proof}

\subsection{Uncontrollable Events and Recovered-State Crashes}
% 不可控事件与恢复状态崩溃
\label{app:crash}

The resolved alphabet contains only commits mediated by the protected monitor.
% 已解析字母表仅包含由受保护监控器中介的提交。
Agent request arrival, scheduling, settlement, delivery, and crash are environmental events, but none appends a \Fresh or \Alias symbol.
% Agent 请求到达、调度、结算、交付与崩溃属于环境事件，但都不会追加 \Fresh 或 \Alias 符号。
Request arrival and scheduler choice have no durable transition until a guarded rule commits.
% 请求到达与调度器选择在某条带守卫规则提交之前没有持久转移。
Settlement and delivery emit modeled API observations but leave the resolved prefix unchanged.
% 结算与交付会发出已建模 API 观测，但保持已解析前缀不变。
\(\Dispatch\) is observable only when the oldest prepared item is durably released, and it also leaves the resolved prefix unchanged.
% \(\Dispatch\) 仅在最早已准备条目被持久释放时可观测，并且它同样保持已解析前缀不变。
Therefore these uncontrollable events stutter with respect to the conditional-nonblocking plant state.
% 因而，这些不可控事件相对于条件非阻塞被控对象状态表现为停顿。
A deployment in which an environmental action itself commits a resolved occurrence is outside this contract and must instead synthesize a supremal controllable sublanguage.
% 如果某个部署中的环境动作本身会提交一个已解析发生项，则它超出本契约范围，必须改为综合上确界可控子语言。

\begin{lemma}[Uncontrollable stutter]
% 不可控事件停顿
\label{lem:app-uncontrollable}
If \(\mathsf{AgentSec}(S;c,r)\), \(e\) is an uncontrollable modeled event, and \(S\xrightarrow eS'\), then the resolved prefix and canonical monitor state are unchanged and \(\mathsf{AgentSec}(S';c,r)\) holds.
% 若 \(\mathsf{AgentSec}(S;c,r)\)、\(e\) 是一个已建模的不可控事件且 \(S\xrightarrow eS'\)，则已解析前缀和规范监控器状态保持不变，并且 \(\mathsf{AgentSec}(S';c,r)\) 成立。
\end{lemma}

\begin{proof}
Request arrival and scheduling have no durable transition, while denial, retry, retrieval, delivery, and crash preserve the recovered security state.
% 请求到达与调度没有持久转移，而拒绝、重试、取回、交付和崩溃保持恢复后的安全状态。
Dispatch and settlement change only monotone outbox, receipt-phase, or result-future fields permitted by execution coherence.
% 派发与结算只改变执行一致性所允许的单调发件箱、回执阶段或结果期值字段。
None changes \(H.T\), \(H.\chi\), the authority projection of \(\Delta\), \(r\), or \(\mathsf q(\eta)\), and the corresponding row of \cref{tab:app-event-matrix} preserves every other invariant clause.
% 这些事件都不改变 \(H.T\)、\(H.\chi\)、\(\Delta\) 的权限投影、\(r\) 或 \(\mathsf q(\eta)\)，并且 \cref{tab:app-event-matrix} 中的对应行保持其他每项不变量条款。
\end{proof}

The crash statement uses the recovered-state interface assumed in \cref{sec:model}, not a new storage protocol.
% 崩溃陈述使用 \cref{sec:model} 假定的恢复状态接口，而不是一种新的存储协议。
Let \(D\) be a concrete durable image, let \(v\) be volatile state, and let \(t\) be at most one in-flight domain transaction.
% 令 \(D\) 为具体持久映像，令 \(v\) 为易失状态，并令 \(t\) 为至多一个进行中的域事务。
Linearizability and crash stability provide a recovery function satisfying
% 可线性化性与崩溃稳定性提供满足下式的恢复函数：
\[
\mathsf{Recover}(D,v,t)=
\begin{cases}
D, & t\text{ has not linearized},\\
\mathsf{commit}_t(D), & t\text{ has linearized}.
\end{cases}
\]
No recovered image contains a strict subset of one modeled atomic update.
% 任何恢复映像都不会包含某项建模原子更新的真子集。
In particular, a \Fresh transaction cannot expose a receipt without its global-trace record, branch-local cursor record, and outbox preparation, and installation cannot activate a new epoch without closing the old epoch and installing all target bindings.
% 特别地，\Fresh 事务不可能在缺少其全局轨迹记录、分支局部游标记录和发件箱准备项时暴露回执，而安装不可能在未关闭旧纪元并安装全部目标绑定时激活新纪元。

\begin{lemma}[Recovered-state crash refinement]
% 恢复状态崩溃精化
\label{lem:app-crash}
If a concrete execution refines compiled state \(S\) immediately before a transaction linearizes, recovery refines \(S\).
% 若具体执行在事务线性化之前立即精化编译状态 \(S\)，则恢复会精化 \(S\)。
If the transaction linearizes to a rule successor \(S'\), recovery refines \(S'\).
% 若事务线性化到规则后继 \(S'\)，则恢复会精化 \(S'\)。
At recovered-state LTS boundaries, crash is consequently a \(\tau\) self-loop.
% 因而，在恢复状态 LTS 边界上，崩溃是一个 \(\tau\) 自环。
\end{lemma}

\begin{proof}
The recovery equation selects the entire preimage or the entire committed postimage.
% 恢复等式会选择完整的提交前映像或完整的已提交后映像。
The former has the same abstraction as \(S\), while the latter has the unique guarded update described by the corresponding row of \cref{tab:app-event-matrix} and therefore abstracts to \(S'\).
% 前者与 \(S\) 具有相同抽象，而后者具有 \cref{tab:app-event-matrix} 对应行所述的唯一带守卫更新，因此抽象为 \(S'\)。
Once the LTS records the recovered image, crashing again changes only volatile state and hence becomes a self-loop.
% 一旦 LTS 记录了恢复映像，再次崩溃只会改变易失状态，因此成为自环。
\end{proof}

\subsection{Durable Agent-API Weak Bisimulation}
% 持久 Agent API 弱互模拟
\label{app:bisimulation}

We now prove the durable-realization clause of \cref{thm:exact-history} by exhausting the observable and silent cases.
% 现在，我们通过穷尽可观测与静默情形来证明 \cref{thm:exact-history} 的持久实现条款。
For \(K=(\kappa,H,\Delta,\mathsf{out},\mathcal A)\), \(\mathsf{Live}(H)=X_{\sem{H.T}}\), and \(I=(K,c_I,r_I)\), let \(\alpha_I(I)=(K,c_I,r_I,H.\eta,H.\beta)\) and \(\alpha_S(S;c,r)=(K,\mathsf{icut}(c),r,H.\eta,\mathsf{bind}|_{\mathsf{Live}(H)})\).
% 对于 \(K=(\kappa,H,\Delta,\mathsf{out},\mathcal A)\)、\(\mathsf{Live}(H)=X_{\sem{H.T}}\) 与 \(I=(K,c_I,r_I)\)，令 \(\alpha_I(I)=(K,c_I,r_I,H.\eta,H.\beta)\)，并令 \(\alpha_S(S;c,r)=(K,\mathsf{icut}(c),r,H.\eta,\mathsf{bind}|_{\mathsf{Live}(H)})\)。
Define \(I\,\mathcal R\,S\) iff some \(c,r\) satisfy \(\mathsf{AgentSec}(S;c,r)\) and \(\alpha_I(I)=\alpha_S(S;c,r)\); write \(\xRightarrow{\ell}=\tau^*\ell\tau^*\) and \(\xRightarrow{\tau}=\tau^*\).
% 定义 \(I\,\mathcal R\,S\) 当且仅当存在 \(c,r\) 满足 \(\mathsf{AgentSec}(S;c,r)\) 与 \(\alpha_I(I)=\alpha_S(S;c,r)\)；记 \(\xRightarrow{\ell}=\tau^*\ell\tau^*\) 且 \(\xRightarrow{\tau}=\tau^*\)。
Fix \(I\,\mathcal R\,S\) with witness \((c,r)\).
% 固定一个以 \((c,r)\) 为见证的 \(I\,\mathcal R\,S\)。
Equality of \(\alpha_I(I)\) and \(\alpha_S(S;c,r)\) gives both machines the same \(\kappa\), current policy, normalized history, receipt ledger, outbox, result futures, semantic cut, resolved prefix, active abstract epoch, and live gate bindings.
% \(\alpha_I(I)\) 与 \(\alpha_S(S;c,r)\) 相等，使两台机器具有相同的 \(\kappa\)、当前策略、规范化历史、回执账本、发件箱、结果期值、语义切点、已解析前缀、活跃抽象纪元和活跃门绑定。

\paragraph{History edges}
% 历史边
The normalization equality makes both machines form the same current \(\Theta_I(u)\) and request-cut token \(\gamma(\Theta_I(u))\).
% 规范化相等使两台机器形成相同的当前 \(\Theta_I(u)\) 与请求切点令牌 \(\gamma(\Theta_I(u))\)。
If derivation is undefined, deterministic failed-check reflection makes both sides emit the same canonical \(\mathsf{Invalid}\) orbit and take a state self-loop.
% 若推导未定义，则确定性的检查失败反射使两侧发出相同的规范 \(\mathsf{Invalid}\) 轨道并执行状态自环。
If derivation exists but no declarative realization exists, \cref{lem:app-spec-bridge} and deterministic fixed-point checking make both sides emit the same canonical \(\mathsf{Reject}\) orbit and take a state self-loop.
% 若推导存在但不存在声明式实现，则 \cref{lem:app-spec-bridge} 与确定性不动点检查使两侧发出相同的规范 \(\mathsf{Reject}\) 轨道并执行状态自环。
Otherwise the ideal machine enables its history edge and, by \cref{lem:app-spec-bridge}, the compiled checker accepts exactly the same history edge and installs the same greatest language.
% 否则，理想机启用其历史边，并且根据 \cref{lem:app-spec-bridge}，编译后检查器恰好接纳同一条历史边并安装相同的最大语言。
The compiled side may take finitely many inactive-preloading \(\tau\)-steps before atomic installation.
% 编译后一侧可以在原子安装之前执行有限次非活跃预加载 \(\tau\) 步。
Pure compilation is outside the recovered-state LTS.
% 纯编译位于恢复状态 LTS 之外。
The ideal side takes its single atomic edge, and both successors normalize to the same \((\kappa_u,H_u,\mathcal A_u)\), cut, empty prefix, active epoch, and target bindings.
% 理想机一侧执行单个原子边，而两个后继都规范化为相同的 \((\kappa_u,H_u,\mathcal A_u)\)、切点、空前缀、活跃纪元和目标绑定。
Deterministic allocation gives the same reply-bundle orbit, so both expose \(\mathsf{Edge}(u,[\mathcal H_u^+]_{\cong})\).
% 确定性分配给出相同的回复包轨道，因此两侧都暴露 \(\mathsf{Edge}(u,[\mathcal H_u^+]_{\cong})\)。
Conversely, every successful compiled installation has passed the schema, fixed-point, seal, and cut checks, so \cref{lem:app-spec-bridge} enables the matching ideal edge.
% 反之，每次成功的编译后安装都已通过模式、不动点、封印和切点检查，因此 \cref{lem:app-spec-bridge} 会启用匹配的理想边。
This argument applies uniformly to all six structural rewrite modes and authenticated live extension.
% 该论证一致地适用于全部六种结构重写模式与经认证的在线扩展。

\paragraph{Protected uses and concrete invocations}
% 受保护使用与具体调用
For a submitted invocation \(m\), both guards compute the same \(m^\circ=\mathsf{canon}_\kappa(m)\).
% 对于提交的调用 \(m\)，两个守卫都会计算相同的 \(m^\circ=\mathsf{canon}_\kappa(m)\)。
The compiled handle check verifies the issuer, scope, occurrence, gate, cell, lineage, signed digest, current epoch, and equality \(m^\circ=\mathsf{inv}_\rho(d)\).
% 编译后句柄检查会验证发行方、作用域、发生项、门、单元、谱系、签名摘要、当前纪元以及等式 \(m^\circ=\mathsf{inv}_\rho(d)\)。
The ideal guard requires the same normalized binding and invocation equality.
% 理想守卫要求相同的规范化绑定与调用相等性。
Both sides then inspect the same receipted-cell set \(D_\Delta\), so they choose the same \(a\in\{\Fresh(x,d),\Alias(x,d)\}\).
% 随后，两侧检查相同的已有回执单元集合 \(D_\Delta\)，因此选择相同的 \(a\in\{\Fresh(x,d),\Alias(x,d)\}\)。
By \cref{cor:app-exact-boundary,lem:compiled-language}, \(ra\) is enabled by the compiled derivative exactly when \(ra\in L^*\) in the ideal machine.
% 根据 \cref{cor:app-exact-boundary,lem:compiled-language}，当且仅当理想机中 \(ra\in L^*\) 时，编译后导数才启用 \(ra\)。

In the \Fresh case, both successors apply the same paired residual, append the same symbol to the global trace and branch-local cursor, advance \(\nu\), create the same cell-bound receipt, and prepare the same canonical invocation.
% 在 \Fresh 情形中，两个后继都应用相同的成对残余、把同一符号追加到全局轨迹与分支局部游标、推进 \(\nu\)、创建同一单元绑定回执，并准备同一规范调用。
The compiled outbox stores \((d,\mathsf{inv}_\rho(d))\), never the caller buffer, while the ideal outbox stores \((d,m^\circ)\); the verified equality makes these records identical.
% 编译后发件箱存储 \((d,\mathsf{inv}_\rho(d))\) 而绝不存储调用方缓冲区，理想发件箱存储 \((d,m^\circ)\)；经验证的等式使这两条记录相同。
In the \Alias case, both successors apply the same paired residual, append the same receipt-silent symbol to the global trace and branch-local cursor, and bind the occurrence to the same durable result future without changing \(\Delta\) or the outbox.
% 在 \Alias 情形中，两个后继都应用相同的成对残余、把同一个不产生回执的符号追加到全局轨迹与分支局部游标，并把该发生项绑定到同一持久结果期值，而不改变 \(\Delta\) 或发件箱。
Both successors remove \(x\)'s live binding; because the residual frontier no longer contains \(x\), their normalized live-binding components remain equal.
% 两个后继都移除 \(x\) 的活跃绑定；由于残余前沿不再包含 \(x\)，它们规范化后的活跃绑定分量仍然相等。
In either case, \cref{lem:app-event-derivative} re-establishes \(\mathcal R\) with witness \((c,ra)\).
% 在任一情形中，\cref{lem:app-event-derivative} 都以 \((c,ra)\) 为见证重新建立 \(\mathcal R\)。

If any occurrence, gate, handle, invocation, epoch, branch, or monitor transition check fails, both normalized guards fail.
% 如果任何发生项、门、句柄、调用、纪元、分支或监控器转移检查失败，则两个规范化守卫都会失败。
Both machines emit \(\mathsf{deny}(x)\), preserve state, and remain related.
% 两台机器都会发出 \(\mathsf{deny}(x)\)、保持状态并维持相关关系。

\paragraph{Use--installation races}
% 使用--安装竞态
Use and installation serialize on the same history and epoch records.
% 使用与安装会在相同的历史和纪元记录上串行化。
If a \Fresh use commits first, it changes \((G,T,\Delta,\chi,\nu,\mathsf{out})\), all of which are bound directly or through \(H\) by the cut seal, so the prepared installation becomes stale.
% 如果 \Fresh 使用先提交，它会改变 \((G,T,\Delta,\chi,\nu,\mathsf{out})\)，这些字段都由切点封印直接绑定或通过 \(H\) 绑定，因此已准备安装会过期。
If an \Alias use commits first, it changes \((G,T,\chi,\nu)\) even though \(\Delta\) and the outbox are unchanged, so the prepared installation again becomes stale.
% 如果 \Alias 使用先提交，即使 \(\Delta\) 和发件箱不变，它仍会改变 \((G,T,\chi,\nu)\)，因此已准备安装同样会过期。
If installation commits first, it atomically closes \(\eta^-\), activates \(\eta^+\), refreshes every live gate and handle, and exposes replacement handles only after activation.
% 如果安装先提交，它会原子关闭 \(\eta^-\)、激活 \(\eta^+\)、刷新每个活跃门与句柄，并且只在激活之后暴露替换句柄。
The old \Fresh or \Alias use then fails the current-epoch guard and is denied without progress.
% 此时，旧的 \Fresh 或 \Alias 使用无法通过当前纪元守卫，并在不产生进度的情况下被拒绝。
The ideal atomic order has exactly the same two cases, so neither side admits a third race outcome.
% 理想原子顺序恰好具有相同的两种情形，因此任一侧都不会接纳第三种竞态结果。

\paragraph{Replies, releases, and silent steps}
% 回复、释放与静默步骤
An authenticated retry, bundle retrieval, or result delivery stutters only in security-state fields and follows the same authenticated link on both sides to expose \(\mathsf{Return}(t,[v]_{\cong})\).
% 经认证重试、包取回或结果交付仅在安全状态字段中表现为停顿，并在两侧沿相同认证链接暴露 \(\mathsf{Return}(t,[v]_{\cong})\)。
\(\Dispatch(d)\) is enabled at the same outbox head and makes the same observable release update.
% \(\Dispatch(d)\) 在相同的发件箱头部启用，并产生相同的可观测释放更新。
\(\mathsf{Settle}(d,[v]_{\cong})\) is enabled for the same receipt and produces the same stable result and normalized update.
% \(\mathsf{Settle}(d,[v]_{\cong})\) 对相同回执启用，并产生相同的稳定结果与规范化更新。
\(\mathsf{Save}\) takes matching \(\tau\)-steps.
% \(\mathsf{Save}\) 执行匹配的 \(\tau\) 步。
Compiled preload and stale or malformed candidate attempts are implementation-only \(\tau\)-steps matched by zero ideal steps because \(\alpha_S\) omits their inactive contents and they emit no verdict.
% 编译后预加载以及过期或畸形候选尝试是仅属于实现的 \(\tau\) 步，由零个理想步骤匹配，因为 \(\alpha_S\) 会忽略其非活跃内容且这些步骤不会发出判定。
Crash is matched by a \(\tau\) self-loop on each recovered-state machine by \cref{lem:app-crash}.
% 根据 \cref{lem:app-crash}，崩溃在每台恢复状态机器上都由一个 \(\tau\) 自环匹配。

\begin{theorem}[Durable Agent-API realization, detailed]
% 持久 Agent API 实现，详细形式
\label{thm:app-bisimulation}
\(\mathcal R\) is a divergence-insensitive weak bisimulation under \(\mathsf{obs}_{\rm api}\).
% \(\mathcal R\) 在 \(\mathsf{obs}_{\rm api}\) 下是发散不敏感弱互模拟。
\end{theorem}

\begin{proof}
The least-generated rule sets in \cref{sec:ideal,sec:kernel} and \cref{tab:app-event-matrix} make the preceding cases exhaustive.
% \cref{sec:ideal,sec:kernel} 中的最小生成规则集与 \cref{tab:app-event-matrix} 使前述情形具有穷尽性。
For every step from either side, the corresponding case constructs a path with observation \(\mathsf{obs}_{\rm api}(\ell)\) and a successor related by the witness in that matrix.
% 对任一侧的每一步，对应情形都会构造一条观测为 \(\mathsf{obs}_{\rm api}(\ell)\) 的路径，以及一个通过该表中见证相关的后继。
The construction is symmetric because successful ideal guards and successful compiled guards were shown equivalent, while implementation-only steps preserve the common abstraction.
% 该构造是对称的，因为成功的理想守卫与成功的编译后守卫已被证明等价，而仅属于实现的步骤保持共同抽象。
Arbitrary repetition of silent preload, stale-candidate, Save, or crash steps is ignored, which is precisely divergence-insensitive weak matching.
% 任意重复的静默预加载、过期候选、Save 或崩溃步骤均被忽略，这恰是发散不敏感弱匹配。
\end{proof}

\subsection{Executable-Artifact Bounds}
% 可执行制品界限

Above 128 outcomes, eight occurrences per outcome, 50,000 linearizations, or 100,000 completions, the executable wrapper returns \(\mathsf{ResourceLimit}(b,n)\) before search and performs no installation.
% 当结果超过 128 个、每个结果的发生项超过 8 个、线性化超过 50,000 个或完成路径超过 100,000 条时，可执行包装器会在搜索前返回 \(\mathsf{ResourceLimit}(b,n)\) 且不执行安装。
This fail-closed engineering outcome is neither a semantic rejection nor an impossibility certificate.
% 这一故障关闭的工程结果既不是语义拒绝，也不是不可能性证书。
The theorem uses the uncapped finite semantic compiler.
% 定理使用无上限的有限语义编译器。

No step in this proof infers a natural-language contract, provides fairness, implements partial fencing inside one domain, or strengthens local outbox release into remote exactly-once execution.
% 本证明中的任何步骤都不推断自然语言契约、不提供公平性、不实现单个域内的部分隔离，也不把本地发件箱释放强化为远程恰好一次执行。
